\section{The motor subsystem}\label{sec:motors}

The motor subsystem is QtRocket's most complete vertical slice: three independent ingest paths
(RockSim XML, RASP text, and the thrustcurve.org REST API) funnel into a single source-agnostic
database, which feeds a thrust-and-mass model that plugs into the composite part tree so that the
rocket's mass, center of gravity, and inertia tensor track the burn (\cref{sec:model}). It is also
the clearest exhibit of the document's recurring motif: the seams are built and tested ahead of
their consumers. The remote-source interface has a fake injected through a deliberate test seam
and the ignition API accepts an arbitrary start time, yet production code ignites exactly one
motor at exactly $t=0$; ejection delays are parsed by every loader, persisted, and consumed by
nothing.

\subsection{Thrust: \code{ThrustCurve}}\label{sec:motors-thrust}

\code{ThrustCurve} stores a vector of (time, thrust) samples in seconds and Newtons, the maximum
sample time, and an ignition-time offset (\src{model/ThrustCurve.h}{12}). \code{getThrust(t)}
subtracts the ignition offset, returns $0\unit{N}$ outside $[0, t_{\max}]$
(\src{model/ThrustCurve.cpp}{50}), returns exact sample values on a hit, and linearly interpolates
inside an interval (\src{model/ThrustCurve.cpp}{47}). A default or empty curve degenerates to a
single $(0,0)$ point so the interval walk always has something to stand on
(\src{model/ThrustCurve.cpp}{17}).

One boundary case earns attention because it used to be a memory-safety bug. Many published curves
begin at $t > 0$ --- the first sample is the first \emph{measured} point, not an origin. When the
query time precedes that first sample, the interval start must be an implicit $(0,0)$ origin:

\begin{listing}[htbp]
\begin{minted}{cpp}
   // Interval start is the previous sample -- unless t precedes the first sample (i == begin), where
   // the curve has no (0,0) origin and we ramp from it. std::prev(begin()) here would read out of
   // bounds.
   double tStart = 0.0;
   double thrustStart = 0.0;
   if(i != thrustCurve.cbegin())
   {
      tStart = std::prev(i)->first;
      thrustStart = std::prev(i)->second;
   }
\end{minted}
\caption{The implicit-origin ramp in \srccap{model/ThrustCurve.cpp}{76}.}
\label{lst:motors-ramp}
\end{listing}

\begin{keyinsight}[Why the origin ramp matters]
Before this guard, the interval walk computed \cppinline|std::prev(i)| with \cppinline|i| at
\cppinline|begin()| --- an out-of-bounds heap read whose value depended on allocation layout, so
early-burn thrust, and therefore apogee, varied from run to run. The fix is pinned by a named
regression test (\src{model/tests/ThrustCurveTests.cpp}{46}), and the interpolation semantics ---
exact samples, midpoints, pre-ignition and post-burnout zeros --- are pinned alongside it
(\src{model/tests/ThrustCurveTests.cpp}{28}). \ok
\end{keyinsight}

\subsection{Mass: the 128-point impulse-proportional curve}\label{sec:motors-mass}

\code{MotorModel} pairs a static \code{MetaData} block --- manufacturer, burn time, impulse class,
delays, propellant and total weight in kilograms (\src{model/MotorModel.h}{285}) --- with runtime
state: \code{ignitionOccurred}, \code{burnOutOccurred}, \code{emptyMass}, \code{isp},
\code{ignitionTime}, and a precomputed \code{massCurve} (\src{model/MotorModel.h}{337}). The
\code{data} member is public with a standing TODO to make it private
(\src{model/MotorModel.h}{282}); it is mutated freely across the codebase in the meantime.

Motor mass over the burn is not taken from data files. Installing metadata triggers
\code{computeMassCurve()} (\src{model/MotorModel.cpp}{85}), which sets
$m_{\text{empty}} = m_{\text{total}} - m_{\text{prop}}$, computes a specific impulse
$I_{\mathrm{sp}} = I_{\text{tot}} / (g_0\, m_{\text{prop}})$ with
$g_0 = 9.80665\unit{m/s^2}$ (\src{utils/math/Constants.h}{10}) --- a value that is then never read
anywhere in the tree --- and forward-Euler-depletes the propellant in proportion to the fraction of
total impulse delivered, i.e.\ a constant-$I_{\mathrm{sp}}$ assumption
(\cref{lst:motors-masscurve}).

\begin{listing}[htbp]
\begin{minted}{cpp}
   // Precompute the mass curve as a lookup table: keeps repeated getMass()/getThrust() calls within
   // one time step consistent, and speeds up the sim. 128 sample points (RASP files cap at 32).
   massCurve.reserve(128);
   double timeStep = data.burnTime / 127.0;
   double t = 0.0;
   double propMass{data.propWeight};
   for(std::size_t i = 0; i < 127; ++i)
   {
      massCurve.push_back(std::make_pair(t + i*timeStep, propMass));
      propMass -= thrust.getThrust(t + i*timeStep) * timeStep * data.propWeight / data.totalImpulse;
   }
   massCurve.push_back(std::make_pair(data.burnTime, 0.0));
\end{minted}
\caption{\code{computeMassCurve()} resamples the thrust integral into a 128-point propellant
lookup table (\srccap{model/MotorModel.cpp}{132}).}
\label{lst:motors-masscurve}
\end{listing}

\code{getMass(simTime)} then has three phases: loaded weight before ignition, empty mass plus
interpolated propellant during the burn, and empty (casing) mass after burnout
(\src{model/MotorModel.cpp}{28}). This is the same fallback a mature tool uses when a file carries
no measured mass data --- but QtRocket applies it \emph{always}, discarding the per-sample mass and
CG measurements the bundled RockSim file actually provides (\cref{sec:motors-ingest}).

\begin{designrule}[Curve before metadata]
Because \code{computeMassCurve()} integrates the thrust curve, the curve must be installed before
the metadata. All three construction sites honor the ordering; it is stated in a comment at two of
them --- ``Order matters'' at \src{model/MotorModelDatabase.cpp}{303} and an equivalent note at
\src{model/ThrustCurveClient.cpp}{328} --- while \src{model/RSEDatabaseLoader.cpp}{96} performs the
same ordering silently. Nothing enforces it: a caller that reverses the two calls gets a mass curve
integrated over an empty curve, with no diagnostic.
\end{designrule}

The computation carries latent hazards, none reachable through current call paths.
\code{computeMassCurve} divides by \code{propWeight}, \code{totalImpulse}, and \code{burnTime}
without zero guards (\src{model/MotorModel.cpp}{135}), so defaulted metadata yields inf/NaN table
entries; \code{massCurve} is reserved and appended but never cleared
(\src{model/MotorModel.cpp}{134}), so a second \code{setMetaData} on the same instance ---
permitted by the public API (\src{model/MotorModel.cpp}{85}), never done by any loader --- would
append a second ramp and break the monotonic interval walk; and \code{getMass} dereferences
\cppinline|massCurve.cbegin()| with no empty check (\src{model/MotorModel.cpp}{40}), so a
default-constructed motor that is started and queried mid-burn is undefined behavior.

\subsection{Ignition and the two burnout clocks}\label{sec:motors-clocks}

The runtime ignition API is one line:
\cppinline|void startMotor(double startTime) { ignitionOccurred = true; ignitionTime = startTime; }|
(\src{model/MotorModel.h}{328}). The parameter suggests staged or delayed ignition; the plumbing
does not deliver it. \code{getMass} honors the offset (it evaluates the curve at
\cppinline|simTime - ignitionTime|, \src{model/MotorModel.cpp}{36}), but \code{getThrust}
passes \emph{absolute} simulation time to the curve (\src{model/MotorModel.cpp}{82}), and the
curve's own offset never moves: \code{ThrustCurve::setIgnitionTime} exists
(\src{model/ThrustCurve.cpp}{41}, with its \code{maxTime} adjustment commented out at line~44) and
has zero call sites in the tree. \code{getThrust} also never consults \code{ignitionOccurred}, so
an un-started motor still reports curve thrust. All of this is masked by a single fact:
\code{RocketModel::launch()} is the only production ignition and always calls
\cppinline|startMotor(0.0)| (\src{model/RocketModel.cpp}{108}). Any future
\code{startMotor}$(t \neq 0)$ would burn propellant on one clock and produce thrust on another.

\begin{hardfail}[Mass and thrust burn out on different clocks]
Even at $t=0$ ignition the two sides use different burnout criteria: \code{getMass} gates the burn
phase on the \emph{metadata} field \code{data.burnTime} (\src{model/MotorModel.cpp}{36}), while
\code{getThrust} gates on the \emph{curve's} last sample time
(\src{model/MotorModel.cpp}{72}). RASP-derived motors are immune because the loader defines
\code{burnTime} as the last sample time, but RockSim files supply the two independently --- the
bundled data carries a \code{burn-time} attribute alongside its sample list
(\src{data/Aerotech.rse}{6}) --- so a file whose attribute disagrees with its samples produces a
window in which thrust is zero while propellant is still depleting, or vice versa. No test covers
the disagreement. \fail
\end{hardfail}

\subsection{Storage and search: \code{MotorModelDatabase}}\label{sec:motors-db}

\code{MotorModelDatabase} is the single ingest and selection surface: the RSE and RASP loaders are
owned and consumed inside it (\src{model/MotorModelDatabase.cpp}{63}), the remote client sits
behind it, and the GUI and CLI only ever call the database. Internally it is a
\cppinline|std::map<std::string, MotorModel>| keyed by motor common name
(\src{model/MotorModelDatabase.h}{129}). The key choice has consequences: same-name motors from
different sources silently replace one another, logged only at debug level
(\src{model/MotorModelDatabase.cpp}{44}) --- two manufacturers who both make a ``D12'' cannot
coexist --- and the import entry points report \emph{net-new} counts, so re-importing a file
reports zero (\src{model/MotorModelDatabase.cpp}{61}, pinned as intentional by
\src{tests/MotorDatabasePersistenceTests.cpp}{348}).

The query side is a small source-agnostic API: \code{MotorQuery}/\code{MotorSummary} and
\code{listMotors} (\src{model/MotorModelDatabase.h}{33}), filtering in memory on exact manufacturer
and impulse class, diameter within $\pm 0.5\unit{mm}$ (a deliberate tolerance for XML float
round-trip fuzz, \src{model/MotorModelDatabase.cpp}{106}), and case-sensitive name substring;
results come back name-sorted for free by map iteration (\src{model/MotorModelDatabase.cpp}{98}).
\code{getMotorModel} returns an \cppinline|std::optional<MotorModel>| by copy
(\src{model/MotorModelDatabase.cpp}{81}), via a \code{find()}-then-\cppinline|operator[]| double
lookup (\src{model/MotorModelDatabase.cpp}{94}).

Persistence is a native XML format (\code{.qmd} by convention): \code{saveMotorDatabase} writes a
\code{QtRocketMotorDatabase} root with a \code{version="0.1"} attribute
(\src{model/MotorModelDatabase.cpp}{184}), one \code{motor} node per entry carrying every
\code{MetaData} field, delays as a CSV string, and the curve as \code{thrust time/force} nodes
(\src{model/MotorModelDatabase.cpp}{177}); the loader mirrors it with per-field defaults, skips
foreign nodes, and recomputes the mass curve on \code{setMetaData}
(\src{model/MotorModelDatabase.cpp}{241}). The round trip over the full bundled import is exact
field-by-field (\src{tests/MotorDatabasePersistenceTests.cpp}{297}), resting on
\code{str()}$\leftrightarrow$\code{toEnum()} inverses pinned for every enum value
(\src{tests/MotorDatabasePersistenceTests.cpp}{245}) --- including, notably, misspellings that
\code{CertOrg::str()} emits (``Austrialian\ldots'', ``Unkown'', \src{model/MotorModel.h}{128}) and
\code{toEnum} explicitly accepts back (\src{model/MotorModel.h}{143}), so the typos are now a
persisted wire format. The \code{version} attribute is written but never examined on load
(\src{model/MotorModelDatabase.cpp}{249}): a future format change has no migration hook.

\subsection{Three ingest paths}\label{sec:motors-ingest}

\begin{description}
\item[RockSim \code{.rse} --- \fpartial.] \code{RSEDatabaseLoader} reads the XML with
  \cppinline|boost::property_tree| and iterates \code{engine-database.engine-list}
  (\src{model/RSEDatabaseLoader.cpp}{24}). It honors average/peak thrust, burn time, code (as the
  common name), the delays CSV, diameter, length, manufacturer, total impulse, and type, converting
  \code{initWt}/\code{propWt} from grams to kilograms (\src{model/RSEDatabaseLoader.cpp}{81});
  impulse class is derived from the code because RSE does not carry it
  (\src{model/RSEDatabaseLoader.cpp}{73}, handling ``1/4A''/``1/2A'' prefixes via
  \src{model/MotorModel.cpp}{97}). What it drops matters more: each \code{eng-data} sample's
  \code{m} (mass) and \code{cg} attributes --- manufacturer-measured mass and CG histories present
  in the bundled file (\src{data/Aerotech.rse}{19}) --- are ignored in favor of the reconstructed
  constant-$I_{\mathrm{sp}}$ curve of \cref{sec:motors-mass}, since the sample loop reads only
  \code{t} and \code{f} (\src{model/RSEDatabaseLoader.cpp}{89}). \code{Isp}, \code{massFrac},
  throat/exit diameters, and the auto-calc flags are likewise ignored; availability is hardcoded
  \code{REGULAR} and the certifying organization \code{UNK}
  (\src{model/RSEDatabaseLoader.cpp}{52}). Robustness is uneven: \code{Type} is the only attribute
  fetched without a default, so a file missing it throws (\src{model/RSEDatabaseLoader.cpp}{85}),
  and delay parsing is unvalidated \code{atoi} --- far laxer than the RASP path below.
\item[RASP \code{.eng} --- \fpresent.] \code{RASPLoader} is a strict, hand-written validating
  parser: a seven-field header checked with typed, line-numbered parse errors
  (\src{model/RASPLoader.cpp}{137}), delays split on \code{-} with \code{P}/\code{none} mapped to
  1000 (the plugged sentinel, \src{model/RASPLoader.cpp}{91}), strictly increasing sample times,
  non-negative thrust, exactly one mandatory trailing zero-thrust sample, and multi-motor files
  (\src{model/RASPLoader.cpp}{265}). It inserts an explicit $(0,0)$ origin when the first sample
  sits at $t>0$ (\src{model/RASPLoader.cpp}{207}) and \emph{derives} burn time, total impulse
  (trapezoidal, \src{model/RASPLoader.cpp}{189}), and average/peak thrust from the samples --- which
  is why RASP motors are immune to the two-clock hazard of \cref{sec:motors-clocks}. RASP weights
  are already kilograms, so no conversion applies (\src{model/RASPLoader.cpp}{153}). Rejection
  behavior is pinned by three negative tests (\src{model/tests/RASPLoaderTests.cpp}{112}, 124,
  138).
\item[thrustcurve.org REST --- \fpartial.] \code{ThrustCurveClient} GETs three unauthenticated v1
  endpoints --- \code{metadata.json} (\src{model/ThrustCurveClient.cpp}{268}),
  \code{search.json} capped at 100 results (\src{model/ThrustCurveClient.cpp}{292}, with truncation
  detected from the server's match count, \src{model/ThrustCurveClient.cpp}{312}), and one
  \code{download.json} per result (\src{model/ThrustCurveClient.cpp}{248}) --- over
  \code{utils::CurlConnection}, a single-handle libcurl GET wrapper with $10\unit{s}$/$30\unit{s}$
  timeouts and default TLS verification (\src{utils/CurlConnection.cpp}{52}). Download parsing
  selects exactly one simfile, preferring the first RASP entry with samples and never concatenating
  alternative series (\src{model/ThrustCurveClient.cpp}{66}); gram-denominated weights are
  converted to kilograms (\src{model/ThrustCurveClient.cpp}{208}). Rough edges: query values are
  concatenated into the URL without percent-encoding (\src{model/ThrustCurveClient.cpp}{285}),
  certifying organization and delays are explicit TODOs (\src{model/ThrustCurveClient.cpp}{197}),
  and a motor whose download returns no samples is still merged into the database carrying the
  default zero-thrust curve (\src{model/ThrustCurveClient.cpp}{324}) --- it would fly with zero
  thrust, a case no flight test exercises. The N+1 downloads run synchronously on the calling
  thread, which in the GUI is the event loop; the interface-level consequences are
  \cref{sec:interfaces}'s subject.
\end{description}

The remote source sits behind the \code{ThrustCurveAPI} interface
(\src{model/ThrustCurveClient.h}{68}), and the database exposes a private constructor taking
\cppinline|std::unique_ptr<ThrustCurveAPI>| to a single test-access friend
(\src{model/MotorModelDatabase.h}{114}), with the real client lazily constructed on first online
use (\src{model/MotorModelDatabase.cpp}{137}). The seam is genuinely consumed: a recording fake
drives the online wrappers offline (\src{tests/MotorDatabasePersistenceTests.cpp}{106}), verifying
criteria mapping --- \code{nameContains} is deliberately not sent upstream
(\src{model/MotorModelDatabase.h}{104}) --- and that merged results become visible to local queries
(\src{tests/MotorDatabasePersistenceTests.cpp}{539}). The pure response parsers are split from the
HTTP layer for the same reason and tested against canned JSON
(\src{model/ThrustCurveClient.h}{88}); the HTTP layer itself is not tested
(\cref{sec:testing}).

\subsection{The motor in flight}\label{sec:motors-flight}

A motor enters the simulation as \code{model::part::Motor}, a leaf \code{Part} owning a
\code{MotorModel} by value (\src{model/parts/Motor.h}{28}). Its sole physics contribution is an
override: \code{getMass(t)} forwards to the motor model (\src{model/parts/Motor.h}{40}), so the
part tree's time-aware composite walk (\cref{sec:model}) produces mass$(t)$, CG$(t)$, and $I(t)$
with no motor-specific inertia code --- the per-unit-mass tensor is a solid cylinder built from the
catalog diameter and length, converted from millimeters at the boundary, with a point-mass fallback
when geometry is unknown (\src{model/parts/Motor.cpp}{27}). The composite behavior is well pinned:
$I(t)$ matches an independent parallel-axis oracle during the burn
(\src{model/tests/MotorTests.cpp}{162}), and post-burnout composites freeze bitwise with no
recompute (\src{model/tests/MotorTests.cpp}{231}). \code{setMotorModel} swaps the model in place,
re-seeding tensor and mass and dirty-flagging ancestors so \code{RocketModel}'s borrowed raw
pointer stays valid (\src{model/parts/Motor.cpp}{16}).

Placement is the weak half. \code{RocketModel::setMotorModel} attaches the single \code{Motor}
child at the airframe root, joining airframe-CM to motor-CM stations with an offset documented as
``Zero today'' (\src{model/RocketModel.cpp}{124}, \src{model/RocketModel.h}{147}): the motor's mass
is coincident with the airframe's center of mass, not seated in an aft motor mount, so the very
CG-shift-during-burn machinery the part tree computes so carefully acts at the wrong station. There
is no user-facing way to reposition it --- a \code{Motor} cannot be added through the part factory,
and a \code{Motor} node inside a design's part tree is rejected as foreign
(\src{model/DesignSerializer.cpp}{137}). The single-motor assumption is explicit:
\code{findMotorInTree} is a DFS for the \emph{first} motor node (\src{model/RocketModel.cpp}{13}),
and \code{Stage.h} is a commented-out ``Not yet'' include (\src{model/RocketModel.h}{23}) --- no
clusters, no stages. In flight, thrust participates as
\cppinline|Vector3{0.0, 0.0, motorPart ? motorPart->getMotorModel().getThrust(t) : 0.0}|
applied through the CM (\src{model/RocketModel.cpp}{70}) with \code{getTorques()} returning zero
unconditionally (\src{model/RocketModel.cpp}{94}); what that simplification costs is
\cref{sec:aero}'s subject. \code{launch()} ignites at exactly $t=0$
(\src{model/RocketModel.cpp}{105}), and \code{isMotorSet()} is the launch gate both front ends read
(\src{model/RocketModel.h}{64}). Designs persist the motor by common name only and re-resolve it
against the database on load, warning rather than failing on a miss
(\src{model/DesignSerializer.cpp}{251}); the format details belong to \cref{sec:persistence}.

Delays and ejection charges deserve their own verdict: \fabsent. Every ingest path parses delay
lists (with 1000 as the plugged sentinel, \src{model/MotorModel.h}{303}), the XML round-trips them,
and no simulation, GUI, or CLI code ever reads them --- there is no ejection event, no recovery
deployment, and the only flight event is descending below the launch altitude
(\src{model/RocketModel.cpp}{62}). Online-sourced motors do not even populate the field
(\src{model/ThrustCurveClient.cpp}{198}). The flight-phase consequences --- no deployment drag
model, descent simulated as ballistic coast --- are taken up in \cref{sec:aero}, and the item is
consolidated with the rest of the gap list in \cref{sec:roadmap}.

\subsection{Bundled data}\label{sec:motors-data}

The repository ships exactly one motor data file, \src{data/Aerotech.rse}{1}, whose 252
\code{engine} entries collapse to 249 stored motors on import: three codes (G77R, G79W, I65W)
appear twice and the name-keyed map of \cref{sec:motors-db} keeps only the last of each, a count
the persistence suite's own setup comment acknowledges
(\src{tests/MotorDatabasePersistenceTests.cpp}{291}). Nothing is loaded at startup:
\code{QTROCKET\_DATA\_DIR} is a tests-only compile definition (\src{tests/CMakeLists.txt}{16}),
so a GUI or CLI user begins with an empty database and must import a file or search online
(\src{gui/CannonballTab.cpp}{137}). A second fixture,
\srcf{tests/data/motors/estes\_small.qmd}, exists purely to exercise the native format. OpenRocket,
by contrast, ships the full thrustcurve.org corpus offline; the breadth comparison is drawn in
\cref{sec:comparison}.
